%%%PAGE 018 rot=0 legibility=good summary=磁场中粒子圆轨迹技巧续：双直线边界（最值相切、时间极值）、反向磁场、隐形磁场、找圆心
%%%BLOCK id=018-01 ch=11 sec=双直线边界·最值相切 kind=derivation alt=none cont=prev y=0.02-0.37
$AD=2R$
%FIG p018_a 0.115 0.045 0.32 0.28
\notefig[0.3]{figures/p018_a.png}
\textbf{最值相切}
%FIG p018_b 0.365 0.10 0.45 0.168
\notefig[0.25]{figures/p018_b.png}
\[ SB^2=R^2-(d-R)^2 \]
\[ SC^2=OT^2=R^2-(d-R)^2 \]
\[ R<d<2R \]
%FIG p018_c 0.68 0.035 0.89 0.285
\notefig[0.3]{figures/p018_c.png}
\[ SB^2=R^2-(R-d)^2 \]
\[ SC^2=R^2-(R-d)^2 \qquad \because SC=SB \]
故\ $SD=2SC$
% 原稿：最后一行"故 SD=2SC"下方有红色下划线
%%%END
%%%BLOCK id=018-02 ch=11 sec=双直线边界·时间 kind=method alt=none cont=none y=0.37-0.60
\textbf{双直线边界}\ $t$\quad 从右边界射出\quad $V$\unclear{同}$\to$弦长最短$\to$圆心角最小
%FIG p018_d 0.12 0.455 0.222 0.595
\notefig[0.25]{figures/p018_d.png}
\[ \sin\theta=\frac{\frac12 d}{R} \]
\[ t_{\min}=\frac{2\theta m}{qB} \]
% 原稿：图右侧另有两条空的竖直长线（未画完的双直线边界图），无其它内容
%%%END
%%%BLOCK id=018-03 ch=11 sec=反向磁场 kind=method alt=none cont=none y=0.60-0.80
\textbf{反向磁场}
%FIG p018_e 0.09 0.637 0.30 0.782
\notefig[0.3]{figures/p018_e.png}
若\ $B_1=B_2$\ $L_1=L_2$\ \underline{则初末速度平行}
% 原稿：此行"则初末速度平行"下方为红色下划线

若\ $B_1\ne B_2$\ $L_1\ne L_2$\ 若初末速度平行\ 则\ \underline{$B_1L_1=B_2L_2$}
% 原稿：此行"B1L1=B2L2"下方为红色下划线，红线向右延伸至页边

速偏角$\alpha_1$，速偏角$\alpha_2$\quad $\alpha_1=\alpha_2\ \overset{\text{相似}}{\Longleftrightarrow}\ \dfrac{R_1}{L_1}=\dfrac{R_2}{L_2}$\quad $R=\dfrac{mv}{qB}$
%%%END
%%%BLOCK id=018-04 ch=11 sec=隐形磁场 kind=method alt=none cont=none y=0.80-0.94
\textbf{隐形磁场}
%FIG p018_f 0.17 0.822 0.42 0.942
\notefig[0.3]{figures/p018_f.png}
\begin{itemize}
\item 最小面积：圆\quad 以入射出射弦为直径的圆
\item 三角形：等边$\triangle$\ 等边三角形
\item 矩形：入射出射弦为长，切弧成形
\item 不定形：粒子未经的地方去除
\end{itemize}
% 原稿：以上四行左侧有大括号；"矩形"一行右端有一个小图（矩形内含一段圆弧）
%FIG p018_g 0.925 0.87 0.995 0.905
\notefig[0.25]{figures/p018_g.png}
%%%END
%%%BLOCK id=018-05 ch=11 sec=找圆心 kind=method alt=none cont=none y=0.94-1.0
\textbf{找圆心}
\begin{itemize}
\item 速度在磁场内：速度作垂线找圆心
\item 速度不在磁场内：速偏角的补角平分线过圆心
\end{itemize}
% 原稿：以上两行左侧有大括号
%%%END
%%%PAGEEND 018
